% ====================================================================
% SECTION 03: ATP STOICHIOMETRY AND CHEMICAL GAUGE INVARIANCE
% ====================================================================
\section{ATP Stoichiometry and Metabolic Scaling Symmetry}
\label{sec:gauge}

\subsection{The Principle of Linear Energy Transduction}

The choice of a cost functional for vascular optimization is not an arbitrary
modeling decision but a consequence of the fundamental \emph{ATP stoichiometry}
of biological transport. In any system operating under metabolic pressure, every
additional Watt of dissipated power $\Delta \Phi$ corresponds to a fixed,
extensive quantity of glucose consumed per unit time. Because glucose is a
fungible good with a constant caloric density, the fitness cost of vascular
inefficiency is strictly proportional to the absolute energetic excess.

This physical constraint imposes a \emph{Metabolic Scaling
Symmetry}\footnote{Not a field-theoretic gauge symmetry (like U(1) or SU(2) in
physics), but a global scaling invariance of the cost functional. In biological
transport, this reflects the fact that fitness selection depends on relative
metabolic excess, not absolute power consumption.}: the fitness penalty must be
invariant under a uniform multiplicative rescaling of the basal metabolism,
$\Phi_{\mathrm{opt}} \to \lambda\Phi_{\mathrm{opt}}$ where $\lambda \in
\mathbb{R}^+$. To satisfy this invariance while maintaining allometric scale
independence (i.e., a dimensionless penalty that allows a shrew and a whale to
be governed by the same adaptive rules), the functional must take the form of a
fractional excess. The linear penalty is thus not a modeling choice, but a
consequence grounded in ATP stoichiometry and constrained by the first law of
thermodynamics.

\subsection{Thermodynamic Linearity and the Physical Occam's Razor}

The linear form of the cost functional $F = (\Phi - \Phi^*)/\Phi^*$ is not an
arbitrary mathematical choice, but a strict consequence of oxidative
phosphorylation stoichiometry and near-equilibrium thermodynamics. At the
cellular level, excess fluidic dissipation is paid in units of ATP. Because this
metabolic accounting is strictly additive---one extra Joule of dissipated
physical work requires a strictly proportional increase in oxidative
metabolism---the evolutionary fitness penalty must scale linearly with excess
entropy production. Imposing non-linear penalties (e.g., quadratic or
logarithmic) would imply unphysical non-stoichiometric metabolic costs,
equivalent to postulating ``compound interest'' on ATP molecules.

Furthermore, numerical sensitivity analysis shows that imposing artificial
non-linear penalties alters the theoretical attractor $\alpha^*$ by less than
$0.3\%$. It is crucial to note that these alternative penalties are merely
counterfactual numerical robustness tests, not physically admissible
alternatives. Their negligible impact is the signature of Prigogine's minimum
entropy production principle in the Onsager regime: because the evolutionary
attractor is strongly confining, the mature vascular network operates in the
immediate vicinity of the fundamental physical lower bound ($\Phi \to \Phi^*$).
In this near-optimal regime, any smooth evolutionary cost landscape is
analytically dominated by its first-order (linear) term.

Finally, the Gauge Invariance (normalization by the intrinsic absolute minimum
$\Phi^*$) represents the allometric ``Occam's razor'': it ensures that the
selective pressure acting on a $30$\,g mouse is mathematically isomorphic to the
pressure acting on a $3000$\,kg elephant, rendering the optimization topology
scale-free and permitting universal scaling laws such as Kleiber's metabolic
rate to emerge naturally.

We formalize this thermodynamic necessity as a first-order response principle:

\begin{theorem}[Onsager Linearity of Fitness Penalty]
\label{thm:onsager}
Consider a vascular network characterized by a state vector $\vec{x} = (r_0,
\ldots, r_G, \alpha, \beta)$ describing vessel radii and branching exponents.
Let $\vec{x}^*$ be the configuration minimizing total dissipation
$\Phi(\vec{x})$ subject to morphometric constraints (fixed total volume,
hierarchical structure, etc.).

For small deviations $\delta\vec{x} = \vec{x} - \vec{x}^*$ from the optimum, the
fitness penalty $\mathcal{L}_{\mathrm{fitness}}$ associated with increased
metabolic cost is:
\begin{equation}
    \mathcal{L}_{\mathrm{fitness}}(\vec{x})
    = \lambda_{\mathrm{bio}} \frac{\Phi(\vec{x}) - \Phi(\vec{x}^*)}{\Phi(\vec{x}^*)}
    + O\left(\|\delta\vec{x}\|^2\right)
\end{equation}
where $\lambda_{\mathrm{bio}} > 0$ is a phenotype-specific selection
coefficient, and the linearity holds to first order in $\|\delta\vec{x}\|$.

Moreover, the functional form $(\Phi - \Phi_{\mathrm{opt}})/\Phi_{\mathrm{opt}}$
is the \emph{uniquely consistent} scale-invariant measure of dissipative excess
satisfying the joint requirements of scale invariance, thermodynamic
extensivity, and normalization.
\end{theorem}

\begin{proof}[Proof sketch; see Supplemental Material, Section S5 for full derivation]
\textbf{Step 1: Near-equilibrium expansion.} For small deviations from the
optimum:
\begin{equation}
    \Phi(\vec{x}^* + \delta\vec{x})
    = \Phi(\vec{x}^*)
    + \frac{1}{2} \delta\vec{x}^T \mathbf{H} \delta\vec{x} + O(\|\delta\vec{x}\|^3),
\end{equation}
where $\mathbf{H}$ is the Hessian matrix (positive semi-definite at a minimum).

\textbf{Step 2: Onsager entropy production.} The excess entropy production rate
is:
\begin{equation}
    \Delta\dot{S} = \frac{\Delta\Phi}{T_{\mathrm{body}}}
    \propto \|\delta\vec{x}\|^2,
\end{equation}
quadratic in deviations (Onsager's theorem for near-equilibrium
systems~\cite{onsager1931}).

\textbf{Step 3: Fractional fitness cost.} Evolutionary selection acts on
fractional metabolic increase, not absolute power (allometric invariance):
\begin{equation}
    \mathcal{L}_{\mathrm{fitness}}
    \sim \frac{\Delta\Phi}{\Phi_{\mathrm{opt}}}.
\end{equation}

The fractional form is necessary because absolute dissipation varies by $10^{5}$
across species (mouse $\sim 0.1$\,W, elephant $\sim 10^4$\,W), but fitness
selection is dimensionless.

\textbf{Step 4: Uniqueness via Compositionality.} The functional $(\Phi -
\Phi_{\mathrm{opt}})/\Phi_{\mathrm{opt}}$ is the simplest form satisfying:
\begin{itemize}
    \item
Scale invariance: $F(\Lambda\Phi, \Lambda\Phi_{\mathrm{opt}}) = F(\Phi,
\Phi_{\mathrm{opt}})$,
    \item
Normalization: $F(\Phi_{\mathrm{opt}}, \Phi_{\mathrm{opt}}) = 0$,
    \item
Monotonicity: $\partial F/\partial\Phi > 0$.
\end{itemize}
While these axioms allow any monotonic $F(x/x_{opt})$, the additional
requirement of \textbf{Compositionality}---that the optimality ranking of
independent subsystems remains invariant under aggregation---singles out the
linear (affine) form.

From scale invariance, $F$ depends only on $\Phi/\Phi_{\mathrm{opt}}$. Write
$F(x) = f(x/1)$ where $x = \Phi/\Phi_{\mathrm{opt}}$. For small $\epsilon = x -
1$:
\begin{equation}
    f(1 + \epsilon) = f(1) + f'(1)\epsilon + O(\epsilon^2) = f'(1)\epsilon,
\end{equation}
using $f(1) = 0$. Setting $f'(1) = 1$ (absorbing into $\lambda_{\mathrm{bio}}$):
\begin{equation}
    \mathcal{C}_{\mathrm{met}} \propto \frac{\Phi - \Phi_{\mathrm{opt}}}{\Phi_{\mathrm{opt}}}.
\end{equation}
\end{proof}

\begin{remark}[Robustness of Linearization in Physiological Regime]
Empirical morphometry (Murray exponent $\alpha = 2.39 \pm 0.15$,
meta-analysis~\cite{taylor2024}) suggests $\|\delta\vec{x}\|/\|\vec{x}^*\|
\lesssim 0.1$, confirming operation in the linear response regime. Explicit
calculation for the human vasculature at $\alpha = 2.39$ yields:
\begin{equation}
    \frac{\Delta\Phi}{\Phi_{\mathrm{opt}}} =
    \frac{\Phi(2.39) - \Phi(2.33)}{\Phi(2.33)} \approx 0.18\text{--}0.22,
\end{equation}
where $\alpha_{\mathrm{opt}} \approx 2.33$ for static transport. This fractional
excess lies well within the regime where second-order corrections $O(\epsilon^2)
\lesssim 0.05$ remain negligible, validating the linear approximation regardless
of the specific functional form $(x-1)$, $(x-1)^2$, or $\ln x$.

While biological systems are fundamentally far-from-equilibrium, the Local
Equilibrium Hypothesis (LEH) is well-justified for vascular remodeling, as the
timescale of angiogenic adaptation is slow relative to the metabolic flux. This
"angiogenic linearity" justifies the quadratic expansion of the excess entropy
production, identifying the selection coefficient as a thermodynamic restoring
force. $\lambda_{\mathrm{bio}} \sim 0.2$--$0.3$ is consistent with
cardiovascular traits under stabilizing selection.
\end{remark}

\subsection{Extensivity and Conditional Uniqueness of the Linear Form}

Before establishing the full scale-invariant functional, we prove that the
additivity (extensivity) of metabolic costs, when assumed, selects the linear
penalty functional as the uniquely consistent form. This \textbf{Canonical
Selection} provides a strong theoretical anchor for the exactness of the
results, though we note that alternative functionals are numerically
near-degenerate: the logarithmic penalty shifts $\alpha^*$ by $\Delta \alpha^*
\approx -0.004$, while the quadratic case yields $\Delta \alpha^* \approx
+0.007$.

\textbf{Biological justification for extensivity.} In evolutionary biology,
fitness costs associated with quantitative traits are commonly assumed to be
additive when measured on fractional or logarithmic scales~\cite{lande1983}. For
metabolic traits, this reflects the stoichiometry of oxidative phosphorylation:
each additional watt of excess dissipation costs a fixed number of ATP
molecules, independent of the baseline metabolic flux. When two independent
vascular networks $A$ and $B$ supply metabolically independent tissues, the
total fitness penalty is the sum of the individual penalties, each normalized by
its respective baseline consumption. Since the total baseline is
$\Phi_{A,\mathrm{opt}} + \Phi_{B,\mathrm{opt}}$, the fractional excess of the
combined system is naturally the weighted average
\begin{equation}
    \mathcal{C}_{A+B} = w_A \mathcal{C}_A + w_B \mathcal{C}_B, \quad
    w_i = \frac{\Phi_{i,\mathrm{opt}}}{\Phi_{A,\mathrm{opt}} + \Phi_{B,\mathrm{opt}}}.
\end{equation}
This \emph{extensibility} principle is an established framework in quantitative
genetics and provides a thermodynamic foundation for the separability condition
invoked below.

\paragraph{Compositionality and Segmentation Invariance.} 
Beyond metabolic stoichiometry, the linear form is uniquely required by the
requirement of \textbf{Segmentation Invariance}. If the vascular network is
partitioned into arbitrary sub-networks (e.g., separating the arterial tree from
the microvascular capillary bed), the total metabolic penalty must be
independent of this arbitrary segmentation. As proven in Supplementary
Section~\ref{S-sec:S7_segmentation}, any nonlinear penalty (e.g., quadratic or
logarithmic) would introduce an artificial scaling dependence on the number of
partitions, making the optimal exponent a function of the observer's chosen
segmentation. The linear scaling form is the unique ``Compositional Point''
where the global optimization is scale-invariant under hierarchical
partitioning.

\begin{theorem}[Conditional Uniqueness of Linear Penalty Given Extensivity]
\label{thm:additivity}
Let $\mathcal{C}(\alpha)$ be a dimensionless penalty functional for the total
metabolic cost, defined on the extensive dissipation power
$\Phi_{\mathrm{net}}(\alpha)$. \emph{If} the following extensivity condition
holds:
\begin{enumerate}
    \item
\textbf{Extensivity}: For two independent vascular subsystems $A$ and $B$ with
dissipations $\Phi_A(\alpha)$ and $\Phi_B(\alpha)$, the total penalty satisfies:
    \begin{equation}
        \mathcal{C}_{A+B} =
        \frac{\Phi_A + \Phi_B - (\Phi_{A,\mathrm{opt}} + \Phi_{B,\mathrm{opt}})}
        {\Phi_{A,\mathrm{opt}} + \Phi_{B,\mathrm{opt}}}
        = \frac{\Phi_{A+B} - \Phi_{A+B,\mathrm{opt}}}{\Phi_{A+B,\mathrm{opt}}}.
    \end{equation}
    \item
\textbf{Separability}: The penalty for the composite system can be expressed in
terms of the individual penalties and their weights:
    \begin{equation}
        \mathcal{C}_{A+B}(\alpha) = w_A \mathcal{C}_A(\alpha) + w_B \mathcal{C}_B(\alpha),
    \end{equation}
where $w_A = \Phi_{A,\mathrm{opt}}/\Phi_{A+B,\mathrm{opt}}$ and $w_B =
\Phi_{B,\mathrm{opt}}/\Phi_{A+B,\mathrm{opt}}$.
\end{enumerate}
Then $\mathcal{C}(\alpha)$ must be the affine (linear) functional:
\begin{equation}
    \mathcal{C}(\alpha) = \frac{\Phi(\alpha) - \Phi_{\mathrm{opt}}}{\Phi_{\mathrm{opt}}}.
    \label{eq:additivity_unique}
\end{equation}
\end{theorem}

\begin{proof}
By the representation lemma (proven in the concluding step), scale invariance
requires $\mathcal{C}(\alpha) = F(\Phi/\Phi_{\mathrm{opt}})$ for some function
$F$ with $F(1) = 0$. Write $x = \Phi/\Phi_{\mathrm{opt}}$.

From condition (2), for subsystems $A$ and $B$:
\begin{equation}
    F\left(\frac{\Phi_A + \Phi_B}{\Phi_{A,\mathrm{opt}} + \Phi_{B,\mathrm{opt}}}\right)
    = w_A F\left(\frac{\Phi_A}{\Phi_{A,\mathrm{opt}}}\right)
    + w_B F\left(\frac{\Phi_B}{\Phi_{B,\mathrm{opt}}}\right).
\end{equation}

Substitute $\Phi_A = x_A \Phi_{A,\mathrm{opt}}$ and $\Phi_B = x_B
\Phi_{B,\mathrm{opt}}$:
\begin{equation}
    F(w_A x_A + w_B x_B) = w_A F(x_A) + w_B F(x_B).
\end{equation}

This is Jensen's functional equation. Assuming $F$ is differentiable,
differentiating both sides with respect to $x_A$ yields:
\begin{equation}
    w_A F'(w_A x_A + w_B x_B) = w_A F'(x_A).
\end{equation}
Dividing by $w_A$ gives $F'(w_A x_A + w_B x_B) = F'(x_A)$. Since this relation
must hold independently of the value of $x_B$, the derivative $F'$ must be a
global constant $c$. Consequently, the second derivative is identically zero
($F''(x) = 0$), forcing $F$ to be strictly affine:
\begin{equation}
    F(x) = c(x - 1),
\end{equation}
where the constant $c$ is absorbed into the selection coefficient. Setting $c =
1$:
\begin{equation}
    \mathcal{C}(\alpha) = \frac{\Phi(\alpha)}{\Phi_{\mathrm{opt}}} - 1
    = \frac{\Phi(\alpha) - \Phi_{\mathrm{opt}}}{\Phi_{\mathrm{opt}}}.
\end{equation}
\end{proof}

\begin{remark}[Stochastic Expected Cost Interpretation]
An alternative justification for the linear form emerges from a stochastic
game-theoretic framework (Paper~II~\cite{paperII}): if the vascular network
faces uncertain metabolic loads with probability distribution $p(\eta)$, then
minimizing the \emph{expected} total cost
$\mathbb{E}_\eta[\mathcal{C}_{\mathrm{transport}} + \eta
\mathcal{C}_{\mathrm{wave}}]$ naturally leads to the linear combination of
penalties. This provides a second, independent pathway to linearity beyond
extensivity.
\end{remark}

\begin{remark}[Physical Interpretation]
The extensivity condition reflects the fundamental additivity of metabolic
costs: if two independent vascular networks each waste a fraction $\epsilon$ of
their baseline metabolism, the combined system also wastes a fraction
$\epsilon$. This is a physical constraint from oxidative phosphorylation
stoichiometry: each additional watt of excess dissipation costs a fixed number
of ATP molecules, independent of the baseline flux.

This extensivity constraint strongly disfavors nonlinear alternatives:
\begin{itemize}
    \item
\textbf{Logarithmic penalty} $\ln(\Phi/\Phi_{\mathrm{opt}})$: Violates
extensivity due to Jensen's Inequality. Because the logarithmic function is
strictly concave, the weighted average of the arguments strictly exceeds the
weighted average of the logarithms: $\ln(w_A x_A + w_B x_B) > w_A \ln x_A + w_B
\ln x_B$. This implies that the total penalty of the composite system grows
superlinearly compared to the sum of its independent parts, introducing an
unphysical compounding effect where larger aggregated systems are penalized
disproportionately more than smaller subsystems. This violates the additive
stoichiometry of ATP consumption, which scales linearly with absolute metabolic
waste regardless of system size.
    \item
\textbf{Quadratic penalty} $(\Phi/\Phi_{\mathrm{opt}} - 1)^2$: Violates
extensivity due to its strict convexity. By Jensen's Inequality, the quadratic
penalty of a composite system with unequal subsystem performance ($\epsilon_A
\neq \epsilon_B$) is always strictly lower than the weighted average of the
individual penalties: $(w_A\epsilon_A + w_B\epsilon_B)^2 < w_A\epsilon_A^2 +
w_B\epsilon_B^2$. This unphysically discounts the evolutionary penalty of
unequal performance, violating the requirement that fitness costs scale
additively with absolute metabolic waste.
\end{itemize}
The linear form emerges as the natural functional compatible with the
thermodynamic requirement that fitness penalties scale with absolute metabolic
waste. While the weighted-average extensivity condition is physically motivated
by Gibbs additivity of free energy rather than rigorously derived from a
variational principle, it provides a strong constraint that singles out the
linear penalty as the simplest consistent form.

\textbf{Physical basis and limitations.} The extensivity condition reflects the
stoichiometric additivity of ATP costs in oxidative phosphorylation and is
consistent with quantitative genetics theory~\cite{lande1983}. However, it is
not mathematically derivable from first principles alone—it represents a
biophysical hypothesis grounded in metabolic biochemistry. The uniqueness result
(Theorem~\ref{thm:additivity}) is therefore \emph{conditional} on this
extensivity assumption. Alternative fitness landscapes (e.g., those involving
resource competition or threshold effects) might violate extensivity, in which
case the logarithmic or quadratic penalties could be biologically relevant.
Empirical robustness tests (Supplemental Material, Section S7) show that
alternative forms shift $\alpha^*$ by less than 0.3\%, suggesting the linear
form is effectively unique within the physiological regime.
\end{remark}

\subsection{Representation Theorem and Gauge Invariance}

\begin{lemma}[Representation Theorem for Gauge-Invariant Functionals]
Let $\mathcal{C}(\alpha)$ be a dimensionless functional of the branching
exponent $\alpha$ defined on the network's total extensive transport power
$\Phi_{\mathrm{net}}(\alpha; \vec{\theta})$, where $\vec{\theta} = \{\mu_f, b,
m_w, Q_0\}$ is the vector of scale and metabolic parameters. If $\mathcal{C}$
satisfies:
\begin{enumerate}
    \item
\textbf{Gauge Invariance:} $\mathcal{C}(\alpha; \Lambda\vec{\theta}) =
\mathcal{C}(\alpha; \vec{\theta})$ for all scale transformations $\Lambda > 0$,
    \item
\textbf{Ground State Normalization:} $\mathcal{C}(\alpha_{\mathrm{opt}}) = 0$,
\end{enumerate}
then $\mathcal{C}$ must have the form:
\begin{equation}
    \mathcal{C}(\alpha) = F\left(\frac{\Phi_{\mathrm{net}}(\alpha)}{\Phi_{\mathrm{net}}(\alpha_{\mathrm{opt}})}\right),
\end{equation}
where $F: \mathbb{R}^+ \to \mathbb{R}$ is an arbitrary function with $F(1) = 0$.
\end{lemma}

\begin{proof}
By Euler's theorem for homogeneous functions, if $\Phi_{\mathrm{net}}(\alpha;
\vec{\theta})$ transforms under global scaling as $\Phi_{\mathrm{net}}(\alpha;
\Lambda\vec{\theta}) = f(\Lambda)\,\Phi_{\mathrm{net}}(\alpha; \vec{\theta})$
for some function $f$, then scale invariance (condition 1) requires:
\begin{equation}
    \mathcal{C}(\alpha; \Lambda\vec{\theta}) =
    \mathcal{C}(\alpha; \vec{\theta}).
\end{equation}

By the Buckingham $\Pi$ theorem, any scale-invariant quantity must be
expressible as a function of dimensionless ratios formed from
$\Phi_{\mathrm{net}}(\alpha)$ and other quantities transforming identically
under $\Lambda$. To avoid introducing arbitrary external parameters, the
reference scale must be intrinsic to the network's phase space. The unique such
reference state is the global transport minimum
$\Phi_{\mathrm{net}}(\alpha_{\mathrm{opt}})$, which also transforms as
$\Phi_{\mathrm{opt}}(\Lambda\vec{\theta}) =
f(\Lambda)\,\Phi_{\mathrm{opt}}(\vec{\theta})$.

Therefore, $\mathcal{C}(\alpha)$ depends only on the dimensionless ratio $x =
\Phi_{\mathrm{net}}(\alpha)/\Phi_{\mathrm{net}}(\alpha_{\mathrm{opt}})$:
\begin{equation}
    \mathcal{C}(\alpha) = F(x), \quad x = \frac{\Phi_{\mathrm{net}}(\alpha)}{\Phi_{\mathrm{net}}(\alpha_{\mathrm{opt}})}.
    \label{eq:ratio_form}
\end{equation}

Condition 2 (ground state normalization) requires $F(1) = 0$, completing the
proof of Eq.~\eqref{eq:ratio_form} and establishing the representation lemma.
\end{proof}

\begin{theorem}[Scale-Free Normalization and Linearity of the Transport Penalty]
\label{thm:gauge}
At fixed body mass $M$ (and thus fixed heart rate $f_H \propto M^{-1/4}$), let
$\Phi_{\mathrm{net}}(\alpha; \vec{\theta})$ be the total extensive transport and
metabolic power dissipated by the network. A valid dimensionless network-level
penalty functional $\mathcal{C}_{\mathrm{transport}}^{\mathrm{net}}(\alpha)$
must satisfy:
\begin{enumerate}
    \item
\textbf{Gauge Invariance (Scale Independence):} $\mathcal{C}(\alpha;
\Lambda\vec{\theta}) = \mathcal{C}(\alpha; \vec{\theta})$.
    \item
\textbf{Positivity and Ground State:} $\mathcal{C}(\alpha) \ge 0$, with
$\mathcal{C}(\alpha_{\mathrm{opt}}) = 0$.
    \item
\textbf{Thermodynamic Linearity (Axiom~3):} The penalty is strictly linear in
the fractional excess dissipation. This linearity is supported by three
independent arguments: (i) extensivity of metabolic costs selects the linear
form as the unique consistent functional (Theorem~\ref{thm:additivity},
conditional on the extensivity hypothesis); (ii) near-equilibrium thermodynamics
(Onsager regime, Theorem~\ref{thm:onsager}); (iii) robustness in the
physiological regime where $|\Delta\Phi/\Phi_{\mathrm{opt}}| \sim 0.2$ makes
higher-order corrections negligible (see Remark following
Theorem~\ref{thm:onsager}).
\end{enumerate}
The canonical functional satisfying these axioms is the fractional excess:
\begin{equation}
    \mathcal{C}_{\mathrm{transport}}^{\mathrm{net}}(\alpha) =
    \frac{\Phi_{\mathrm{net}}(\alpha) - \Phi_{\mathrm{net}}(\alpha_{\mathrm{opt}})}
    {\Phi_{\mathrm{net}}(\alpha_{\mathrm{opt}})}.
    \label{eq:gauge_invariant_cost}
\end{equation}
\end{theorem}


\begin{proof}
\textbf{Independence of the three axioms.} Axioms~1 and~2 together permit any
positive-definite function of the dimensionless ratio $x =
\Phi_{\mathrm{net}}(\alpha)/\Phi_{\mathrm{net}}(\alpha_{\mathrm{opt}})$
satisfying $F(1)=0$, such as $(x-1)^2$ or $\ln x$. Axiom~3 is therefore
logically independent and strictly necessary to isolate the affine form.

\textbf{Justification of linearity.} The linear penalty is not axiomatically
imposed but follows from either:
\begin{enumerate}[label=(\alph*)]
    \item
\textbf{Extensivity} (Theorem~\ref{thm:additivity}): if assumed, uniquely
selects linear form;
    \item
\textbf{Onsager near-equilibrium regime} (Theorem~\ref{thm:onsager}):
empirically verified for vascular remodeling timescales;
    \item
\textbf{Physiological near-optimality}: $|\Delta\alpha/\alpha| \sim 0.1$ ensures
higher-order corrections $O(\epsilon^2) \lesssim 0.05$ remain negligible.
\end{enumerate}
Since (b) and (c) are empirical facts, linearity is \emph{conditionally
necessary} rather than axiomatically imposed.

\textbf{Derivation.} By Euler's theorem for homogeneous functions and the
Buckingham $\Pi$ theorem, Axiom~1 requires $\mathcal{C}(\alpha) =
F(\Phi_{\mathrm{net}}(\alpha)/\Phi_{\mathrm{ref}})$, where $\Phi_{\mathrm{ref}}$
transforms identically under $\Lambda$. To avoid introducing arbitrary external
parameters, $\Phi_{\mathrm{ref}}$ must be an intrinsic property of the network's
phase space. The unique such reference state is the global transport minimum
$\Phi_{\mathrm{net}}(\alpha_{\mathrm{opt}})$.

By the representation lemma, $\mathcal{C}(\alpha) = F(x)$ with $x =
\Phi/\Phi_{\mathrm{opt}}$. Axiom~2 gives $F(1) = 0$ and $F(x) \ge 0$ for $x \ge
1$.

To determine the unique form of $F$, write $x = 1 +
\Delta\Phi/\Phi_{\mathrm{opt}}$. Axiom~3 (thermodynamic linearity) requires that
the penalty be strictly linear in the extensive excess entropy production
$\Delta\Phi$, meaning $\partial\mathcal{C}/\partial\Phi$ is constant. This
forces $F$ to be affine in $x$:
\begin{equation}
    F(x) = a(x - 1) + b, \quad a, b \in \mathbb{R}.
\end{equation}

Applying the ground state normalization $F(1) = 0$ gives $b = 0$. Requiring
$F(x) > 0$ for $x > 1$ (positivity of penalty for excess dissipation) gives $a >
0$. Up to a positive multiplicative constant, the unique solution is:
\begin{equation}
    F(x) = x - 1,
\end{equation}
yielding Eq.~\eqref{eq:gauge_invariant_cost}.

\textbf{Consistency established:} The fractional excess $(\Phi_{\mathrm{net}} -
\Phi_{\mathrm{opt}})/\Phi_{\mathrm{opt}}$ is the \emph{uniquely consistent}
functional satisfying the thermodynamic requirements of extensivity and
scale-invariant selection.

\textbf{Exclusion of the logarithm.} The logarithmic map
$\ln(\Phi/\Phi_{\mathrm{opt}})$, for instance, introduces a concave profile that
dampens the evolutionary penalty for large structural deviations, decoupling it
from the linear stoichiometry of oxidative phosphorylation.
\end{proof}

\begin{remark}[Epistemological Status: Parsimony Assumption vs. First Principle]
\label{rem:gauge_parsimony}
Metabolic duty cycle invariance---requiring the cost functional to be unchanged
under $\vec{\theta} \to \Lambda \vec{\theta}$---is not a first principle but an
\textbf{assumption of parsimony} justified by allometric universality. Empirical
metabolic rate scales as $M^{3/4}$ across $10^5$-fold mass
range~\cite{savage2004}. Any violation (i.e., penalty functional depending on
absolute $\theta$) would require a biologically privileged reference scale
(e.g., $M_{\mathrm{human}} = 70\,$kg) that is absent from the data. We therefore
impose scale-freeness as the simplest hypothesis consistent with observation.
Alternative formulations introducing absolute thresholds
$\theta_0(\mathrm{taxon})$ could be tested if future data reveal systematic
species-specific deviations from $M^{3/4}$ scaling; current evidence shows no
such structure.
\end{remark}

\begin{remark}[Heart Rate as External Control Parameter]
The heart rate $f_H$ is \emph{not} part of the scaling group $\vec{\theta}$
because it enters the Womersley number $\mathrm{Wo} = r\sqrt{2\pi f_H/\nu}$
through the pulsatile frequency, not as a metabolic cost factor. The allometric
scaling $f_H \propto M^{-1/4}$ acts as an \emph{external control parameter} that
drives the system through the viscous-to-wave transition as body mass varies.
This explains why $\alpha^*$ is universal \emph{within} a given mass class
(e.g., all 70 kg mammals) but varies \emph{between} mass classes (shrew vs.
elephant). The scale invariance ensures universality at fixed $M$; the
allometric transition arises from the $M$-dependence of $f_H$.
\end{remark}

\begin{remark}
Theorem~\ref{thm:gauge} establishes that
$\mathcal{C}_{\mathrm{transport}}^{\mathrm{net}}$ is not a phenomenological
ansatz but the \emph{only} dimensionless functional consistent with the symmetry
of the biological optimization problem. The ground state $\Phi_{\mathrm{opt}}$
plays the role of a natural gauge: it is intrinsic to the network, transforms
homogeneously with $\vec{\theta}$, and produces the unique invariant ratio.
Paper~III~\cite{paperIII} rigorously derives the metabolic scaling exponent
$\beta(\alpha,d) = d\alpha/(2d+\alpha)$ from proximal dominance (Theorem~1),
showing that scale invariance uniquely maps local branching geometry to global
organismal scaling laws.
\end{remark}

\begin{remark}[Gauge Invariance vs.\ Dimensional Homogeneity]
\label{rem:gauge_vs_dimensional}
It is crucial to distinguish \textbf{Scale-Free Normalization} (Axiom~1) from
simple dimensional homogeneity. Dimensional analysis (Buckingham $\Pi$ theorem)
merely requires that the functional depend on dimensionless combinations of
variables---but it does \emph{not} uniquely specify which combinations or how
they enter.

Scale invariance is a \emph{physical symmetry principle}: the cost functional
must remain invariant under global rescaling transformations $\vec{\theta} \to
\Lambda\vec{\theta}$ that leave the physiological state unchanged (e.g.,
doubling all metabolic rates while doubling all flow rates). This symmetry
\emph{constrains the functional form}: it forces the cost to depend only on
ratios $\Phi(\alpha)/\Phi_{\mathrm{opt}}$, not on their absolute magnitudes or
arbitrary external scales.

Combined with the linearity axiom (Axiom~3), scale invariance uniquely
determines the cost functional to be the fractional excess $(\Phi -
\Phi_{\mathrm{opt}})/\Phi_{\mathrm{opt}}$, excluding all other dimensionless
forms (e.g., logarithmic, quadratic). This is strictly stronger than dimensional
analysis, which would permit any function $F(\Phi/\Phi_{\mathrm{opt}})$.
\end{remark}

\begin{remark}[Gauge Invariance vs.\ Field Theory Gauge]
\label{rem:gauge_terminology}
The term ``scale invariance'' here refers to the freedom to choose the reference
scale $\Phi_{\mathrm{opt}}$ as the metabolic duty cycle—analogous to fixing the
zero of the electrostatic potential in classical electrodynamics, not to local
gauge symmetries (U(1), SU(2)) of field theory. The invariance
$\mathcal{C}(\Lambda\vec{\theta}) = \mathcal{C}(\vec{\theta})$ is a \emph{global
scaling symmetry} that, combined with normalization $F(1)=0$, uniquely
determines the functional form—precisely the role of gauge-fixing in physics.
This terminology emphasizes that the cost functional depends only on
dimensionless ratios, not on arbitrary external scales.
\end{remark}

\begin{remark}[Non-substitutability and aggregation robustness]
The linear combination in $\mathcal{L}_{\mathrm{net}}$ is the unique aggregation
operator consistent with the physical non-substitutability of the two cost
modes: wave energy loss and metabolic excess are orthogonal failure modes that
cannot compensate each other. A multiplicative combination $\mathcal{L} \propto
\mathcal{C}_{\mathrm{wave}}^{\mathrm{net}} \cdot
\mathcal{C}_{\mathrm{transport}}^{\mathrm{net}}$ would permit arbitrarily large
wave dissipation to be offset by arbitrarily efficient transport---a
biologically inadmissible trade-off in a system where signal propagation and
metabolic supply are independently required for viability. The minimax of a
linear combination is the canonical robust-optimization formulation for
non-fungible constraints~\cite{ben-tal2009}, and the $L_q$ robustness analysis
of Paper~II~\cite{paperII} confirms that the saddle point $\alpha^* =
\VarAlphaStar$ is stable across all convex aggregations $q \geq 1$, with a total
spread $\Delta\alpha^* = 0.008$.
\end{remark}

\begin{remark}[Third Commensurable Costs]
One might ask whether the inclusion of a third physiological cost could render
the optimization commensurable. The answer depends strictly on its physical
dimensions. If the third cost is commensurable with metabolism (e.g., another
energetic loss measured in watts), it is simply absorbed linearly into the
transport gauge $\Phi_{\mathrm{net}}$, leaving the dimensionless structure of
the Lagrangian unchanged. If it is incommensurable, it cannot linearly rescue
the local optimization; rather, it introduces a third orthogonal axis to the
minimax phase space, requiring a generalized zero-sum game over three competing
invariants.
\end{remark}

\subsection{Architectural Decoupling: A Consistency Check}

The scale invariance established above has a direct mathematical consequence:
the minimax branching exponent $\alpha^*$ is determined entirely by
dimensionless structural quantities and is \emph{independent of every absolute
metabolic parameter}. This architectural decoupling serves as a consistency
check on the theory.

\begin{corollary}[Structural Ground State and Scale-Free Normalization]
\label{thm:rigidity}
Let $\alpha^*$ be the minimax saddle point of the network Lagrangian
$\mathcal{L}_{\mathrm{net}}(\alpha,\eta)$ for an idealized symmetric,
non-tapering vascular tree. Then $\alpha^*$ depends primarily on the
dimensionless structural parameters $\mathcal{S} = (G, N, p, \alpha_w, A)$ and
exhibits the following scaling behavior:
\begin{equation}
    \alpha^* = \alpha^*_{\mathrm{ground}}(\mathcal{S}) + \delta\alpha^*_{\mathrm{pert}}(\mu_f, \rho, \dots)
    \label{eq:ground_state}
\end{equation}

\textbf{Pure metabolic parameters} ($b$, $Q_0$, $\ell_0$, $\Delta
G_{\mathrm{ATP}}$) exhibit scale invariance: $|S_{\theta_i}| < 0.01$
(Table~\ref{tab:sensitivity}), confirming that the architecture is decoupled
from absolute metabolic scales.

\textbf{Fluid-mechanical parameters} ($\mu_f$, $\rho$) show moderate sensitivity
($|S_x| \approx 0.15$--$0.20$), reflecting their role in defining boundary
conditions and the viscous-wave transition regime. The predicted ground state
$\alpha^*_{\mathrm{ground}} \approx \VarAlphaStarModel$ represents the idealized
symmetric attractor; structural heterogeneities (asymmetry, taper, elastic wall
compliance) introduce perturbations $\delta\alpha^* \approx +\VarAlphaResidual$,
shifting large-mammal values toward the observed $\alpha^* \approx
\VarAlphaStar$.
\end{corollary}

\begin{proof}
By Theorem~\ref{thm:gauge}, the transport cost
$\mathcal{C}_{\mathrm{transport}}^{\mathrm{net}}(\alpha)$ is degree-zero
homogeneous in pure metabolic parameters $\vec{\theta}_{\mathrm{met}}$, which
cancel identically in the ratio $(\Phi_{\mathrm{net}} -
\Phi_{\mathrm{opt}})/\Phi_{\mathrm{opt}}$. The wave cost
$\mathcal{C}_{\mathrm{wave}}^{\mathrm{net}}(\alpha)$ depends on the Womersley
number $\mathrm{Wo} \propto r \sqrt{2\pi f_H / \nu}$ (where $\nu = \mu_f/\rho$
is kinematic viscosity) and geometric branching ratios determined by
$\mathcal{S}$. Changes in viscosity shift the viscous-wave transition boundary,
introducing moderate sensitivity to the minimax balance point. The ground state
$\alpha^*_{\mathrm{ground}}(\mathcal{S})$ is the solution for the idealized
symmetric model; real biological systems exhibit perturbative shifts due to
morphometric heterogeneity and fluid-mechanical variations.
\end{proof}

\begin{remark}[Numerical Verification]
Table~\ref{tab:sensitivity} confirms the two-regime structure: pure metabolic
parameters show $|S_x| < 0.01$ (scale invariance), while fluid-mechanical and
structural parameters show moderate sensitivity ($0.04 \leq |S_x| \leq 1.2$)
across all mammalian species, confirming the robustness of the minimax
attractor.
\end{remark}

\begin{remark}[Physical Interpretation]
The Structural Ground State theorem establishes that the branching exponent is
an \emph{architectural attractor} rather than a rigidly fixed eigenvalue. The
ground state $\alpha^*_{\mathrm{ground}} \approx \VarAlphaStarModel$ emerges
from the network topology and dimensionless structural parameters, while
moderate perturbations from fluid-mechanical variations and morphometric
heterogeneities shift real systems toward $\alpha^* \approx \VarAlphaStar$. This
two-scale structure---a robust topological attractor with physiological
perturbations---explains why the same approximate value appears across species
whose metabolic rates and cardiac outputs differ by orders of magnitude, while
still allowing for the subtle allometric trends observed in empirical data. The
architecture is not \emph{tuned} to the physiology; it is \emph{attracted} to a
structural ground state and modulated by biological perturbations.
\end{remark}
